% Part:sets-functions-relations
% Chapter: size-of-sets
% Section: non-enumerability-alt

\documentclass[../../../include/open-logic-section]{subfiles}

\begin{document}

\olfileid[ar]{sfr}{siz}{nen-alt}

\olsection{مجموعات \printtoken{S}{nonenumerable}}

\begin{editorial}
نبين هنا أن $\Bin^\omega$ و$\Pow{\Nat}$ غير قابلتين للتعداد، ونعتمد تعريفات \olref[enm-alt]{sec}. فالمشترط دالة تقابلية مع~$\Nat$، لا دالة شاملة من~$\PosInt$.
\end{editorial}

\begin{explain}
مجموعة الأعداد الطبيعية~$\Nat$ لا نهائية، وكونها !!{enumerable} بين بنفسه. وإنما الأمر اللافت وجود مجموعات \emph{!!{nonenumerable}}، أي ليست !!{enumerable} (انظر \olref[sfr][siz][enm-alt]{defn:enumerable}).

وقد يستغرب هذا؛ فإن كون $\OLId{latin-looped}{A}$ !!{nonenumerable} يقتضي \emph{انتفاء} كل !!{bijection} $\OLId{latin-ordinary}{f}
\colon \Nat \to \OLId{latin-looped}{A}$. بل لا توجد دالة ترسل !!{element}s~$\Nat$، على كثرتها التي لا تنتهي، إلى~$\OLId{latin-looped}{A}$ وتستوفي~$\OLId{latin-looped}{A}$. فإذا كانت $\OLId{latin-looped}{A}$ !!{nonenumerable}، ففي~$\OLId{latin-looped}{A}$ !!{element}s «أكثر» من الأعداد الطبيعية.

فسبيل إثبات أن المجموعة !!{nonenumerable} أن ننفي وجود !!{bijection} على الوجه المطلوب. ويكفي لذلك أن نبرهن أن كل محاولة لتعداد !!{element}s~$\OLId{latin-looped}{A}$ تفوت !!{element} واحدًا على الأقل؛ فتكون كل دالة $\OLId{latin-ordinary}{f}\colon \Nat \to \OLId{latin-looped}{A}$ غير !!{surjective}. وطريقة كانتور \emph{القطرية} سبيل عام إلى هذا النفي: تعطى قائمة من !!{element}s~$\OLId{latin-looped}{A}$، ولتكن $\OLId{latin-ordinary}{x}_\OLMathNumeral{1}$، $\OLId{latin-ordinary}{x}_\OLMathNumeral{2}$، \dots، فننشئ !!{element} من~$\OLId{latin-looped}{A}$ يمتنع، بحكم إنشائه، أن يكون في القائمة.

وليتضح الأمر بالمثال، نبدأ بالمجموعة~$\Bin^\omega$ التي تضم جميع السلاسل اللانهائية من~$\OLMathNumeral{0}$ و~$\OLMathNumeral{1}$. (فالرمز «$\Bin$» يدل على الثنائية، ويكفينا أن نعده اسمًا للمجموعة ذات العنصرين $\{\OLMathNumeral{0},\OLMathNumeral{1}\}$.)\oliflabeldef{sfr:card-arithmetic:card-opps:sec}{\footnote{والتحقيق أن $\Bin^\omega$ مجموعة جميع متتاليات $\omega$ المؤلفة من~$\OLMathNumeral{0}$ و~$\OLMathNumeral{1}$، أي المجموعة $\funfromto{\omega}{\{\OLMathNumeral{0},\OLMathNumeral{1}\}}$. وإنما يتبين معنى هذا في \olref[sfr][card-arithmetic][card-opps]{sec}.}{} أما هذه العبارة التي فيها شيء من التجوز فلا يترتب عليها لبس فيما نحن فيه.}
\end{explain}

\begin{thm}
\ollabel{thm:nonenum-bin-omega}
$\Bin^\omega$~!!{nonenumerable}.
\end{thm}

\begin{proof}
خذ تعدادًا كيف اتفق لمجموعة جزئية من~$\Bin^\omega$. فتكون بين أيدينا قائمة $\OLId{latin-ordinary}{s}_{\OLMathNumeral{0}}$، $\OLId{latin-ordinary}{s}_{\OLMathNumeral{1}}$، $\OLId{latin-ordinary}{s}_{\OLMathNumeral{2}}$، \dots{}، وكل~$\OLId{latin-ordinary}{s}_\OLId{latin-ordinary}{n}$ سلسلة لا نهائية من~$\OLMathNumeral{0}$ و~$\OLMathNumeral{1}$. وللرمز $\OLId{latin-ordinary}{s}_\OLId{latin-ordinary}{n}(\OLId{latin-ordinary}{m})$ مؤشران: $\OLId{latin-ordinary}{m}$ يحدد رتبة الرقم، و$\OLId{latin-ordinary}{n}$ يحدد السلسلة. فنرسم القائمة مصفوفة يقع فيها $\OLId{latin-ordinary}{s}_\OLId{latin-ordinary}{n}(\OLId{latin-ordinary}{m})$ في الصف ذي الرتبة~$\OLId{latin-ordinary}{n}$ والعمود ذي الرتبة~$\OLId{latin-ordinary}{m}$:

\[
\begin{array}{c|c|c|c|c|c}
& \OLMathNumeral{0} & \OLMathNumeral{1} & \OLMathNumeral{2} & \OLMathNumeral{3} & \dots \\\hline
\OLMathNumeral{0} & \mathbf{s_{0}(0)} & \OLId{latin-ordinary}{s}_{\OLMathNumeral{0}}(\OLMathNumeral{1}) & \OLId{latin-ordinary}{s}_{\OLMathNumeral{0}}(\OLMathNumeral{2}) & \OLId{latin-ordinary}{s}_\OLMathNumeral{0}(\OLMathNumeral{3}) & \dots \\\hline
\OLMathNumeral{1} & \OLId{latin-ordinary}{s}_{\OLMathNumeral{1}}(\OLMathNumeral{0})& \mathbf{s_{1}(1)} & \OLId{latin-ordinary}{s}_\OLMathNumeral{1}(\OLMathNumeral{2}) & \OLId{latin-ordinary}{s}_\OLMathNumeral{1}(\OLMathNumeral{3}) & \dots \\\hline
\OLMathNumeral{2} & \OLId{latin-ordinary}{s}_{\OLMathNumeral{2}}(\OLMathNumeral{0})& \OLId{latin-ordinary}{s}_{\OLMathNumeral{2}}(\OLMathNumeral{1}) & \mathbf{s_2(2)} & \OLId{latin-ordinary}{s}_\OLMathNumeral{2}(\OLMathNumeral{3}) & \dots \\\hline
\OLMathNumeral{3} & \OLId{latin-ordinary}{s}_{\OLMathNumeral{3}}(\OLMathNumeral{0})& \OLId{latin-ordinary}{s}_{\OLMathNumeral{3}}(\OLMathNumeral{1}) & \OLId{latin-ordinary}{s}_\OLMathNumeral{3}(\OLMathNumeral{2}) & \mathbf{s_3(3)} & \dots \\\hline
\vdots & \vdots & \vdots & \vdots & \vdots & \mathbf{\ddots}
\end{array}
\]
وننشئ سلسلة لا نهائية~$\OLId{latin-ordinary}{d}$ من~$\OLMathNumeral{0}$ و~$\OLMathNumeral{1}$ لا ترد في هذه القائمة. وطريق إنشائها تعيين خاناتها كلها، أي تعيين $\OLId{latin-ordinary}{d}(\OLId{latin-ordinary}{n})$ لكل $\OLId{latin-ordinary}{n} \in
\Nat$. فنقرأ قطر المصفوفة---ومن هنا اسم الطريقة القطرية---ونستبدل كل~$\OLMathNumeral{1}$ بـ~$\OLMathNumeral{0}$، وكل~$\OLMathNumeral{0}$ بـ~$\OLMathNumeral{1}$. وضبط ذلك أن نجعل $\OLId{latin-ordinary}{d}(\OLId{latin-ordinary}{n})$ مساويًا لـ$\OLMathNumeral{0}$ أو~$\OLMathNumeral{1}$ تبعًا لكون !!{element} ذي الرتبة~$\OLId{latin-ordinary}{n}$ على القطر، أعني $\OLId{latin-ordinary}{s}_\OLId{latin-ordinary}{n}(\OLId{latin-ordinary}{n})$، مساويًا لـ$\OLMathNumeral{1}$ أو~$\OLMathNumeral{0}$:

\[
\OLId{latin-ordinary}{d}(\OLId{latin-ordinary}{n}) =
\begin{cases}
\OLMathNumeral{1} & \text{إذا كان $s_{n}(n) = 0$}\\
\OLMathNumeral{0} & \text{إذا كان $s_{n}(n) = 1$}
\end{cases}
\]
فيتحقق $\OLId{latin-ordinary}{d} \in \Bin^\omega$؛ إذ هي سلسلة لا نهائية من~$\OLMathNumeral{0}$ و~$\OLMathNumeral{1}$. وقد عينّا~$\OLId{latin-ordinary}{d}$ بحيث $\OLId{latin-ordinary}{d}(\OLId{latin-ordinary}{n}) \neq \OLId{latin-ordinary}{s}_\OLId{latin-ordinary}{n}(\OLId{latin-ordinary}{n})$ لكل $\OLId{latin-ordinary}{n} \in
\Nat$؛ فهي، أي~$\OLId{latin-ordinary}{d}$، تخالف~$\OLId{latin-ordinary}{s}_\OLId{latin-ordinary}{n}$ في الخانة ذات الرتبة~$\OLId{latin-ordinary}{n}$. إذن $\OLId{latin-ordinary}{d} \neq \OLId{latin-ordinary}{s}_\OLId{latin-ordinary}{n}$ لكل $\OLId{latin-ordinary}{n}\in \Nat$، فلا تكون~$\OLId{latin-ordinary}{d}$ في القائمة $\OLId{latin-ordinary}{s}_\OLMathNumeral{0}$، $\OLId{latin-ordinary}{s}_\OLMathNumeral{1}$، $\OLId{latin-ordinary}{s}_\OLMathNumeral{2}$، \dots

فمهما أعطينا تعدادًا لمجموعة جزئية من $\Bin^\omega$، وجدنا !!{element} من~$\Bin^\omega$ خارجًا عنه. فلا تعداد يستوفي المجموعة~$\Bin^\omega$؛ أي إن~$\Bin^\omega$ !!{nonenumerable}.
\end{proof}

\begin{editorial}
\textbf{تنبيه تصحيحي على الأصل (OLSIZ-008).}
عكس الأصل وصفي المؤشرين؛ فالجدول يضع $\OLId{latin-ordinary}{s}_\OLId{latin-ordinary}{n}(\OLId{latin-ordinary}{m})$ في الصف $\OLId{latin-ordinary}{n}$
والعمود $\OLId{latin-ordinary}{m}$، ولذلك هي الخانة $\OLId{latin-ordinary}{m}$ من السلسلة $\OLId{latin-ordinary}{n}$.

\textbf{تنبيه تصحيحي على الأصل (OLSIZ-009).}
كرر الأصل التحويل من $\OLMathNumeral{1}$ إلى $\OLMathNumeral{0}$ وأسقط التحويل المقابل؛
فأُثبت التحويلان $\OLMathNumeral{1}\mapsto\OLMathNumeral{0}$ و$\OLMathNumeral{0}\mapsto\OLMathNumeral{1}$ اللذان تعطيهما صيغة الحالات.
\end{editorial}

\begin{explain}
واسم «الاستدلال القطري» مأخوذ من استعمال قطر المصفوفة في تعريف~$\OLId{latin-ordinary}{d}$. غير أن المصفوفة ليست من لوازمه؛ فقد نثبت أن مجموعة !!{nonenumerable} بمعنى مماثل، من غير مصفوفة مرسومة ولا قطر ظاهر. وبيان ذلك في النتيجة الآتية.
\end{explain}

\begin{thm}
\ollabel{thm:nonenum-pownat}
$\Pow{\Nat}$ ليست !!{enumerable}.
\end{thm}

\begin{proof}
نسلك المسلك نفسه، فنبين أن أي قائمة من المجموعات الجزئية من~$\Nat$ تفوتها مجموعة جزئية من~$\Nat$. ولتكن $\OLId{latin-looped}{N}_\OLMathNumeral{0}, \OLId{latin-looped}{N}_\OLMathNumeral{1}, \OLId{latin-looped}{N}_\OLMathNumeral{2}, \ldots$ قائمة مفروضة من المجموعات الجزئية من~$\Nat$. ونعين مجموعة~$\OLId{latin-looped}{D}$ بهذه الصفة: $\OLId{latin-ordinary}{n} \in \OLId{latin-looped}{D}$ إذا وفقط إذا كان $\OLId{latin-ordinary}{n} \notin
\OLId{latin-looped}{N}_{\OLId{latin-ordinary}{n}}$، أي
\[
\OLId{latin-looped}{D} = \Setabs{\OLId{latin-ordinary}{n} \in \Nat}{\OLId{latin-ordinary}{n} \notin \OLId{latin-looped}{N}_\OLId{latin-ordinary}{n}}
\]
فإن $\OLId{latin-looped}{D}\subseteq \Nat$ ظاهر. وأما~$\OLId{latin-looped}{D}$ فلا ترد في القائمة؛ إذ $\OLId{latin-ordinary}{n} \in \OLId{latin-looped}{D}$ يكافئ $\OLId{latin-ordinary}{n}\notin \OLId{latin-looped}{N}_\OLId{latin-ordinary}{n}$ بمقتضى التعريف، ومنه $\OLId{latin-looped}{D} \neq \OLId{latin-looped}{N}_\OLId{latin-ordinary}{n}$ لكل $\OLId{latin-ordinary}{n} \in \Nat$.
\end{proof}

\begin{explain}
لم نذكر قطرًا في البرهان، ولكن يظهر معناه إذا رسمنا المجموعات $\OLId{latin-looped}{N}_\OLMathNumeral{0}$، $\OLId{latin-looped}{N}_\OLMathNumeral{1}$، \dots{} في مصفوفة. فنجعل $\OLId{latin-looped}{N}_\OLId{latin-ordinary}{n}$ في الصف ذي الرتبة~$\OLId{latin-ordinary}{n}$، ونكتب~$\OLId{latin-ordinary}{m}$ في العمود ذي الرتبة~$\OLId{latin-ordinary}{m}$ إذا وفقط إذا كان $\OLId{latin-ordinary}{m} \in \OLId{latin-looped}{N}_\OLId{latin-ordinary}{n}$. ولو كانت المجموعات الأربع الأولى $\{\OLMathNumeral{0},\OLMathNumeral{1},\OLMathNumeral{2},\dots\}$، $\{\OLMathNumeral{1}, \OLMathNumeral{3}, \OLMathNumeral{5}, \dots\}$، $\{\OLMathNumeral{0},\OLMathNumeral{1},\OLMathNumeral{4}\}$، و$\{\OLMathNumeral{2},\OLMathNumeral{3},\OLMathNumeral{4},\dots\}$، لابتدأت المصفوفة هكذا:
\[
\begin{array}{r@{}rrrrrrr}
  \OLId{latin-looped}{N}_\OLMathNumeral{0} = \{ & \mathbf{0}, & \OLMathNumeral{1}, & \OLMathNumeral{2}, & & & & \dots\}\\
  \OLId{latin-looped}{N}_\OLMathNumeral{1} = \{ &  & \mathbf{1}, &  & \OLMathNumeral{3}, &  & \OLMathNumeral{5}, & \dots\}\\
  \OLId{latin-looped}{N}_\OLMathNumeral{2} = \{ & \OLMathNumeral{0}, & \OLMathNumeral{1}, &  &  & \OLMathNumeral{4}\phantom{,} &  & \}\\
  \OLId{latin-looped}{N}_\OLMathNumeral{3} = \{ &  &  & \OLMathNumeral{2}, & \mathbf{3}, & \OLMathNumeral{4}, & & \dots\}\\
  &\vdots & & & & & & \ddots\phantom{\}}
  \end{array}
\]
فنحصل على~$\OLId{latin-looped}{D}$ بالنزول على القطر، جاعلين $\OLId{latin-ordinary}{n} \in \OLId{latin-looped}{D}$ إذا وفقط إذا كان~$\OLId{latin-ordinary}{n}$ \emph{غائبًا} عن موضعه القطري. ففي هذا المثال نستبعد~$\OLMathNumeral{0}$ و~$\OLMathNumeral{1}$، ونثبت~$\OLMathNumeral{2}$، ونستبعد~$\OLMathNumeral{3}$، ثم نمضي على ذلك.
\end{explain}

\begin{prob}
أقم حجة قطرية صريحة تبين أن مجموعة جميع الدوال $\OLId{latin-ordinary}{f} \colon \Nat \to
\Nat$ !!{nonenumerable}. أي بين أن كل قائمة $\OLId{latin-ordinary}{f}_\OLMathNumeral{1}$، $\OLId{latin-ordinary}{f}_\OLMathNumeral{2}$، \dots{} من الدوال التي نوع كل منها $\OLId{latin-ordinary}{f}_\OLId{latin-ordinary}{i}\colon \Nat \to \Nat$ تفوتها دالة $\OLId{latin-ordinary}{g} \colon \Nat \to \Nat$ من النوع نفسه.
\end{prob}

\end{document}
